% Part: first-order-logic
% Chapter: models-theories
% Section: size-of-structures

\documentclass[../../../include/open-logic-section]{subfiles}

\begin{document}

\olfileid{fol}{mat}{siz}

\olsection{Nyatakake Ukuran \printtoken{P}{structure}}

\begin{explain}
Ana sawenehing sipat struktur kang bisa kita nyatakake
sanajan tanpa nggunakake simbol-simbol nonlogis sawijining
basa. Upamane, ana !!{sentence}s kang bener ing
!!a{structure} yen lan mung yen !!{domain} saka
!!{structure} kasebut ngemot paling ora, paling akeh,
utawa persis sawenehing cacah~$n$ !!{element}s.
\end{explain}

\begin{prop}
!!{sentence}
\begin{multline*}
  !A_{\ge n} \ident \lexists[x_1][\lexists[x_2][\dots\lexists[x_n][{}]]]\\
  \begin{aligned}
  (\eq/[x_1][x_2] \land {}
  \eq/[x_1][x_3] \land \eq/[x_1][x_4] \land \dots \land \eq/[x_1][x_n] \land {}\\
\eq/[x_2][x_3] \land \eq/[x_2][x_4] \land \dots \land {} \eq/[x_2][x_n] \land {} \\
\vdots\\
\eq/[x_{n-1}][x_n])
\end{aligned}
\end{multline*}
bener ing !!a{structure}~$\Struct M$ yen lan
mung yen $\Domain M$ ngemot paling ora
$n$ !!{element}s. Mula,
$\Sat{M}{\lnot !A_{\ge n+1}}$ yen lan mung
yen $\Domain M$ ngemot paling akeh~$n$
!!{element}s.

\end{prop}

\begin{prop}
!!{sentence}
\begin{multline*}
!A_{= n} \ident \lexists[x_1][\lexists[x_2][\dots\lexists[x_n][{}]]] \\
\begin{aligned}
  (\eq/[x_1][x_2] \land {}
  \eq/[x_1][x_3] \land \eq/[x_1][x_4] \land \dots \land \eq/[x_1][x_n] \land {}\\
\eq/[x_2][x_3] \land \eq/[x_2][x_4] \land \dots \land {} \eq/[x_2][x_n] \land {} \\
\vdots\\
\eq/[x_{n-1}][x_n] \land {} \\
\lforall[y][(\eq[y][x_1] \lor \dots \lor \eq[y][x_n]]))
\end{aligned}
\end{multline*}
bener ing !!a{structure}~$\Struct M$ yen lan
mung yen $\Domain M$ ngemot persis
$n$ !!{element}s.
Cathetan notasi: rong tampilan iki pola skematis
kanggo cacah unsur umum; kasus cilik kang ora
nduweni pasangan variabel beda perlu diwaca
kanthi konjungsi kosong kang trep, dudu
kanthi nganggep indeks ing pola mau tetep ana. Nalika cacah siji, konjungsi pasangan beda iku bener kanthi vakum nanging syarat kabeh unsur padha karo paseksi tunggal tetep perlu. Nalika cacah nol, disjungsi kosong iku luput lan domain kosong ora diidini ing definisi struktur iki.
\end{prop}

\begin{prop}
!!{structure} iku tanpa wates yen lan mung
yen dadi model saka
\[
\{!A_{\ge 1}, !A_{\ge 2}, !A_{\ge 3}, \dots \}.
\]
\end{prop}

Ora ana siji ukara kang mung ngemot simbol logis
kang bener ing~$\Struct M$ yen lan mung yen
$\Domain M$ tanpa wates. Nanging, kita bisa
menehi !!{sentence}s kanthi !!{predicate}s
nonlogis kang mung nduweni model tanpa wates
(sanajan ora saben !!{structure} tanpa wates
dadi modele). Sipat dadi struktur winates, lan
sipat dadi struktur !!{nonenumerable}, malah
ora bisa dinyatakake nganggo himpunan
!!{sentence}s tanpa wates. Kasunyatan-kasunyatan
iki dadi akibat saka teorema kompaktitas lan
L\"owenheim--Skolem.

\end{document}
